\section{Conclusion}
\label{sec:conclusion}

We presented the Engram Hybrid Adaptive Consensus, a CometBFT-based protocol
that folds peripheral network health -- Bitcoin settlement, Celestia data
availability, and P2P quality -- into the consensus object as a first-class,
quorum-agreed variable, degrading through a four-state finite state machine
into a local fallback that keeps producing blocks without external
corroboration instead of halting when a peripheral
layer fails. We backed this design with two complementary forms of evidence.
Formally, each protocol layer's own bound-model-checked safety and liveness
properties found zero violations across 3.39 million distinct explored
states ($1.02\times10^{9}$ generated) at that layer's own checked-in scale, and a five-mechanism ablation
study at the FSM layer showed that removing any one safety mechanism
reproduces a distinct, predicted failure mode, evidence that each mechanism
is individually load-bearing at that layer; the cross-layer refinement
check that would extend this guarantee to the full composed system at
production scale is specified but not yet run to completion, and we regard
the live evaluation below as currently carrying the larger share of
whole-system evidence. Empirically, on a real 7-validator testbed (with two
exceptions run at 4-node scale, Section~\ref{sec:setup}) against live
Bitcoin and Celestia networks, the protocol (i) maintained \textbf{safety}
with zero AppHash divergence under representative Byzantine attacks, while
improving \textbf{liveness} by sustaining continuous block production and
completing a full recovery cycle under a layered composite fault; (ii)
provided \textbf{controlled recovery}, with hysteresis measurably trading
uptime against flapping resistance in a tunable way across all three
hysteresis-gated edges; and (iii) incurred \textbf{acceptable overhead} --
+228 bytes per block and $\sim$18.25\,$\mu$s of validation cost on the
healthy path, with no throughput ceiling found within the tested range (up
to 100 tx/s). Together, these results indicate that treating
peripheral-layer health as a first-class consensus variable is a practical
way to trade a small, bounded overhead for substantially stronger liveness
guarantees in modular blockchain systems, though the whole-system claim
currently rests more heavily on live evidence than on the formal side,
for reasons detailed in Section~\ref{sec:discussion}.

\subsection*{Future Work}

Eight directions follow directly from the limitations discussed in
Section~\ref{sec:discussion}. First, the cross-layer refinement
verification sweep (Section~\ref{sec:formal}) remains an open problem
rather than a scheduled task: even its smallest configuration (C1, $N{=}4$)
has not reached convergence, and given the combinatorial state-space growth
typical of explicit-state model checking over composed protocol layers, we
cannot commit to a timeline or a resource budget that would close it.
Narrowing this gap may call for a different verification strategy
altogether -- symmetry reduction, statistical model checking, or a
deductive proof in the style discussed in Section~\ref{sec:related} --
rather than simply more compute directed at the same explicit-state search.
Second, re-running the live evaluation across genuinely distributed
physical hosts, rather than one shared machine, would remove the
single-host resource-contention confound noted in
Section~\ref{subsec:limitations}. Third, closing the gap between the target
DA sampling design and the current single-bridge-per-validator prototype
(Section~\ref{subsec:sensors}) would strengthen the DA-outage results from
measuring a bridge failure to measuring a genuinely sampling-detected
partial DA failure. Fourth, the open question raised in
Section~\ref{sec:discussion} -- whether the
\textsc{Anchored}$\leftrightarrow$\textsc{Suspicious} boundary needs the
same exponential-backoff hardening already applied to the
\textsc{Sovereign}/\textsc{Recovering} edges -- remains an open design
problem for future protocol iterations. Fifth, re-running the
malicious-proposer test and the A1/A2 eclipse/Sybil scenarios at the same
7-node scale used elsewhere would remove the scale inconsistency noted in
Section~\ref{sec:setup}. Sixth, testing $f$ simultaneously colluding
Byzantine validators ($f{=}2$ at $N{=}7$) rather than one actor at a time,
and specifically testing a false individual health-certificate submission
into the vote-extension aggregation, would close the gap between the
formally-bounded threat model and what has actually been exercised live
(Section~\ref{subsec:limitations}). Seventh, repeating the composite-fault,
seven-scenario, P4, and reorg-depth experiments across multiple independent
trials -- as already done for the hysteresis sweep and the overhead
benchmark -- would upgrade their single-trial timings to a statistically
characterized result. Eighth, and most directly actionable, running a
persistent 3-of-7 or 4-of-7 split-view scenario -- where validators'
external reads of Bitcoin, DA, or P2P health genuinely disagree near the
quorum boundary, rather than the decisive 1-of-7 and 2-of-7 splits already
tested -- would close the near-quorum limitation noted in
Section~\ref{subsec:limitations} and clarify whether it explains the
round-skip overhead observed even on the healthy RQ4 path.

\subsection*{Artifact Availability}

% TODO: dien thong tin that -- vd repo public/private, giay phep, va co
% san khi nao (upon publication / ngay bay gio). Day la 1 paragraph chuan
% trong cac paper he thong/security, reviewer se tim muc nay.
The implementation, TLA+ specifications, and experiment scripts used to
produce the results in this paper are available at
\texttt{[TODO: repository URL]}.